% output=pdftex

% jafoe foXet

\usemodule[abr-01]

\definetypesetting[fo][--foxet][frame=on]

\usetypescript[palatino][texnansi]
\setupbodyfont[palatino,10pt]

% \definecolor[red]  [darkred]
% \definecolor[green][darkgreen]
% \definecolor[blue] [darkblue]

\setuplayout
  [backspace=2cm,topspace=2cm,
   width=middle,height=middle]

\setupwhitespace
  [big]

\setupcolors
  [state=start]

\setuphead
  [chapter]
  [style=\bfc]

\setuphead
  [section]
  [style=\bfb]

\setuphead
  [subsection]
  [style=\bfa]

\definefont[BigFont][RegularBold at 75pt]
\definefont[MedFont][RegularBold at 45pt]

% \definecolor[pagebgcolor][s=.85]
% \definecolor[pagefgcolor][s=.15]

\definecolor[pagebgcolor][s=.15]
\definecolor[pagefgcolor][s=.85]

\startuseMPgraphic{page}
  StartPage ;
  picture p, q ;

  numeric alternative ; alternative = 2 ;

  fill Page withcolor \MPcolor{pagebgcolor} ;

  if alternative = 1 :
    p := textext("\tttf XMLFOE") rotated 90 ;
  elseif alternative = 2 :
    p := textext("\chardef\foxetcolormode=0 \foXet") rotated 90 ;
  elseif alternative = 3 :
    p := textext("\tttf foXet") rotated 90 ;
  fi ;
  p := p shifted -llcorner p ;
  p := p xysized(bbwidth(Page),bbheight(Page)) ;
  q := p reflectedabout(llcorner p,ulcorner p) shifted (bbwidth(p),0);
  draw p withcolor transparent(1,0.25,red) ;
  draw q withcolor transparent(1,0.25,green) ;
  p := p rotatedaround(center p, 180) ;
  q := p reflectedabout(llcorner p,ulcorner p) shifted (bbwidth(p),0);
  draw p withcolor transparent(1,0.25,blue) ;
  draw q withcolor transparent(1,0.25,yellow) ;

  if alternative = 1 :
    p := textext("\tttf XMLFOE") ;
  elseif alternative = 2 :
    p := textext("\chardef\foxetcolormode=0 \foXet") ;
  elseif alternative = 3 :
    p := textext("\tttf foXet");
  fi ;
  p := p shifted -llcorner p ;
  p := p xysized(bbwidth(Page),bbheight(Page)) ;
  q := p reflectedabout(llcorner p,lrcorner p) shifted (bbheight(p),0);
  draw p withcolor transparent(1,0.25,red) ;
  draw q withcolor transparent(1,0.25,green) ;
  p := p rotatedaround(center p, 180) ;
  q := p reflectedabout(llcorner p,urcorner p) shifted (bbheight(p),0);
  draw p withcolor transparent(1,0.25,blue) ;
  draw q withcolor transparent(1,0.25,yellow) ;

  StopPage ;
  setbounds currentpicture to Page ;
\stopuseMPgraphic

\defineoverlay[page][\useMPgraphic{page}]

\starttext

\setupbackgrounds[page][background=page]

\startstandardmakeup[headerstate=none,footerstate=none]
  \hfill \BigFont \pagefgcolor\setstrut \strut just an
  \vskip0pt
  \hfill \BigFont \pagefgcolor \setstrut \strut XML FO
  \vskip-15pt
  \hfill \BigFont \pagefgcolor \setstrut \strut engine
  \vfill
  \hfill \MedFont \pagefgcolor Hans Hagen
  \vskip 30pt
  \hfill \MedFont \pagefgcolor PRAGMA ADE
\stopstandardmakeup

% \setuphead[chapter][numbercommand=\FOxETc]
% \setuphead[section][numbercommand=\FOxETs]

\setupbackgrounds[page][background=]

% \noheaderandfooterlines Coming soon! \stoptext

\chapter{Introduction}

\section{Implementation}

When reading the following sections, you need to be aware of the fact
that \FOXET\ is capable of more that just processing formatting
objects. Actually, in many cases, directly mapping \XML\ structure
onto \CONTEXT\ feature is more efficient, and if you want to do more
complex things, it may be a far better alternative.

As a (potential) user you probably want to know how conformant the
implementation is to the specs. Our aim is to implement most of the
features that make sense for printing and screen based interactive
documents.

We will not discuss all features in detail since there are enough
documents around that sort of explain formatting objects. As usual the
formal specification leaves some room for interpretation. While
writing \FOXET\ I was reminded of the process of implementing \MATHML:
using specs that were written with no implementation around to show
the intentions. This also automatically leads to the feeling that we
may need to provide alternative solutions and that is something we
will certainly do.

When looking around for additional resources I noticed that these are
mostly written from the perspective of programming than from the point
of view of design and high end typesetting. In addition to the
specification I used Neil Bradley's XSL companion (first print). I
also occasionally looked into Sebastian Rahtz list of attributes
supported by his \type {passivetex}'s in order to get an idea of what
needed to be supported.

\section{Conformance}

I will not fall into the trap of filling hundreds of pages with boring
lists of attributes but for the moment stick to the remark that most
of what makes sense is somehow provided. You should keep in mind that
formatting objects are somewhat related to cascading style sheets and
so to web publishing and not so much to typography. Formatting objects
describe the page in a sequence of content blocks, fonts and spacing
on a sequence of pages. How exactly things end up on paper is left to
a typesetting engine, and the result therefore depends on the
algorithms used to typeset text and characteristics of fonts.

Although formatting objects are supposedly not to rely too much on any
intelligent from the side of a typesetting engine, there are a few
areas where it has to leave things to the engine, for instance in
handling footnotes and floats. When you are familiar with the kind of
typesetting that I have to deal with in my daily work, you can imagine
a few more aspects of going from \XML\ to print that are outside the
scope of formatting objects. But the majority of documents don't need
anything fancy and presumably can be handles by formatting objects,
given that the attributes are used properly. As said, many of those
document can probably be handled more efficiently (and with better
results) by direct mappings, if only because in that case the
typesetting engine knows what it is dealing with and can act properly
in border cases.

One of the aspects of conformance is the ability to deal with
shortcuts. For some obscure reason the rather verbose formatting
objects source files may be (neclectably) reduced in size by packing
attributes in shortcuts. Apart from the fact that this posses some
limitations on extensibility (due to clashing keywords in such a
packed sequence) it complicates an implementation. One is supposed to
generate those files by means of for instance \XSLT, so there is no
real urging need to pack attributes this way.

Much of the work in writing \FOXET\ went into deciphering the specs,
getting an idea of what was needed when and keying all those
properties.  Most of the machinery needed to implement the specs was
already present in the underlying typesetting engine \CONTEXT\
although often with a different usage in mind. In due time I may
replace some of the current mappings to simplified versions of the
underlying mechanism which may speed up processing. On the other hand,
using more clever mechanisms below the surface gives room for
experiments and extensions.

(Remark: the inheritance model is rather messy.)

\section {Extensions}

The core functionality is provided in the \type {fo} namespace.

\starttabulate
\NC fo \NC basic formatting objects    \NC \NR
\NC fe \NC basic formatting extensions \NC \NR
\NC fx \NC extra formatting objects    \NC \NR
\NC ft \NC user  formatting objects    \NC \NR
\NC fd \NC setup (definitions)         \NC \NR
\stoptabulate


\section {MathML}

\section {Graphics}

\section {Processing}

...


% \stoptext


\chapter{Declarations and pagination and layout formatting objects}

\section{Introduction}

\section{fo:root} % 6.4.2

This object is provides the normal document wrapper. The \type {media-usage}
attribute is not supported since it does not make sense for paper documents.

\showXMLfile{fo-0101.fo}

\section{fo:declarations} % 6.4.3

This object encapsulates definitions meant for the engine. Although it makes
more sense to do this in configuration files, the following is possible:

\showXMLfile{fo-0102.fo}

Those familiar with \CONTEXT\ may recognize these definitions. And yes, you
may specify colors the way you like: \RGB, \CMYK, spotcolors, transparencies,
and of course grayscales.

\section{fo:color-profile} % 6.4.4

\showXMLfile{fo-0103.fo}

\section{fo:page-sequence} % 6.4.5

The page sequence object holds together (part of) the content, for instance
a chapter. The way the pages in this sequence are set up is determined by
the masters.

\showXMLfile{fo-0201.fo}

\section{fo:layout-master-set} % 6.4.6

This is just a wrapper for the page dimensions and sequence
specification.

\section{fo:page-sequence-master} % 6.4.7
\section{fo:single-page-master-reference} % 6.4.8
\section{fo:repeatable-page-master-reference} % 6.4.9
\section{fo:repeatable-page-master-alternatives} % 6.4.10
\section{fo:conditional-page-master-reference} % 6.4.11
\section{fo:simple-page-master} % 6.4.12

\section{fo:region-body} % 6.4.13

The main document content ends up in this region, which may spread
over multiple pages.

\showXMLfile{fo-0613.fo}

Normally a text in multiple columns will have small margins. \in
{Figure} [fig:columns] shows a few examples.

\placefigure
  [here]
  [fig:columns]
  {Using the \type {precedence} attribute in regions.}
  {\startcombination[3*1]
     {\typesetfile[fo][fo-0611.fo][width=.2\textwidth]} {no columns}
     {\typesetfile[fo][fo-0612.fo][width=.2\textwidth]} {two column}
     {\typesetfile[fo][fo-0613.fo][width=.2\textwidth]} {three columns}
   \stopcombination}

You can add graphics to a region. \in {Figure} [fig:backgroundimage]
demonstrates how you can position the graphic.

\placefigure
  [here]
  [fig:backgroundimage]
  {Using the \type {background-image} attribute in regions.}
  {\startcombination[4*1]
     {\typesetfile[fo][fo-0641.fo][width=.2\textwidth]} {repeat}
     {\typesetfile[fo][fo-0642.fo][width=.2\textwidth]} {no-repeat/percentage}
     {\typesetfile[fo][fo-0643.fo][width=.2\textwidth]} {no-repeat/keyword}
     {\typesetfile[fo][fo-0643.fo][width=.2\textwidth]} {no-repeat/dimension}
   \stopcombination}

\section{fo:region-before} % 6.4.14

This region determines what in most cases is considered to be the
header. This area can have backgrounds and frames. The main thing here
is not to mix up the  page and region dimensions and to keep in mind
that the body sort of overlaps with  the regions, unless the margins
are set large enough.

\showXMLfile{fo-0601.fo}

The objects \type {region-before} and \type {region-after} may extend
into the margins depending on the \type {precedence} attribute.

\showXMLfile{fo-0604.fo}

\in {Figure} [fig:precedence] demonstrates the consequences of setting the
\type {precedence} attribute.

\placefigure
  [here]
  [fig:precedence]
  {Using the \type {precedence} attribute in regions.}
  {\startcombination[4*1]
     {\typesetfile[fo][fo-0601.fo][width=.2\textwidth]} {false false}
     {\typesetfile[fo][fo-0602.fo][width=.2\textwidth]} {false true}
     {\typesetfile[fo][fo-0603.fo][width=.2\textwidth]} {true  false}
     {\typesetfile[fo][fo-0604.fo][width=.2\textwidth]} {true  true}
   \stopcombination}

\section{fo:region-after} % 6.4.15

This region is the counterpart of the previously discussed one, and is
often called the footer. For more details, see \type {fo:region-before}.

\section{fo:region-start} % 6.4.16

This region is also known as the left margin (in left to right typesetting).
It has the same properties as \type {fo:region-before}.

\section{fo:region-end} % 6.4.17

This region is also known as the right margin (in left to right typesetting).
It has the same properties as \type {fo:region-before}.

\section{fo:flow} % 6.4.18

\section{fo:static-context} % 6.4.19

This is just a container for content that will end up in a region, for
instance in \type {xsl-region-before}.

\section{fo:title} % 6.4.20

This object does not make sense in paper documents. If users see an application
for it, I will implement a suitable visualization.

\chapter{Block-level formatting objects}

\section{Introduction} % 6.5.1
\section{fo:block} % 6.5.2
\section{fo:block-container} % 6.5.3

\chapter{Inline-level formatting objects}

\section{Introduction} % 6.6.1
\section{fo:bidi-override} % 6.6.2
\section{fo:character} % 6.6.3

The character object can be used for special purposes. You should keep
in mind that the settings are local (which is good).

\showXMLfile{fo-0301.fo}

\in {Figure} [fig:character] demonstrates the results. Forget about the
extra \type {foe} attributes. We need them to get a nice picture.

\placefigure
  [here]
  [fig:character]
  {Some examples of playing with character properties.}
  {\typesetfile[fo][fo-0301.fo][width=.8\textwidth]}

\section{fo:initial-property-set} % 6.6.4
\section{fo:external-graphic} % 6.6.5
\section{fo:instream-foreign-object} % 6.6.6
\section{fo:inline} % 6.6.7
\section{fo:inline-container} % 6.6.8
\section{fo:leader} % 6.6.9
\section{fo:page-number} % 6.6.10

\showXMLfile{fo-0801.fo}

\placefigure
  [here]
  [fig:pagenumber]
  {Some examples of page number formatting and referencing.}
  {\startcombination[2*1]
     {\typesetfile[fo][fo-0601.fo][page=1,width=.4\textwidth]} {Page 123}
     {\typesetfile[fo][fo-0602.fo][page=2,width=.4\textwidth]} {Page 124}
   \stopcombination}

\section{fo:page-number-citation} % 6.6.11

See previous section.

\chapter{Formatting objects for tables}

\section{Introduction} % 6.7.1
\section{fo:table-and-caption} % 6.7.2
\section{fo:table} % 6.7.3
\section{fo:table-column} % 6.7.4
\section{fo:table-caption} % 6.7.5
\section{fo:table-header} % 6.7.6
\section{fo:table-footer} % 6.7.7
\section{fo:table-body} % 6.7.8
\section{fo:table-row} % 6.7.9
\section{fo:table-cell} % 6.7.10

\chapter{Formatting objects for lists}

\section{Introduction} % 6.8.1
\section{fo:list-block} % 6.8.2
\section{fo:list-item} % 6.8.3
\section{fo:list-item-body} % 6.8.4
\section{fo:list-item-label} % 6.8.5

\chapter{Dynamic effects: link and multi formatting objects}

\section{Introduction} % 6.9.1
\section{fo:basic-link} % 6.9.2
\section{fo:multi-switch} % 6.9.3
\section{fo:multi-case} % 6.9.4
\section{fo:multi-toggle} % 6.9.5
\section{fo:multi-properties} % 6.9.6
\section{fo:multi-property-set} % 6.9.7

\chapter{Out-of-line formatting objects}

\section{Introduction} % 6.10.1
\section{fo:float} % 6.10.2
\section{fo:footnote} % 6.10.3
\section{fo:footnote-body} % 6.11.4

\chapter{Other formatting objects}

\section{Introduction} % 6.11.1
\section{fo:wrapper} % 6.11.2
\section{fo:marker} % 6.11.3
\section{fo:retrieve-marker} % 6.11.4

\section{Remarks}

For the moment we don't support functions. Actually, they are a kind
of  contradiction to the fact that formatting objects are sort of
expanded  and don't have any intelligence built in. Functions like
\type  {label-end()} could as well have been replaced by a keyword.

\section {Extensions}

\stoptext